\section{Method}
\label{sec:method}

We study how the choice of interior collocation points affects the accuracy and
convergence of physics-informed neural networks (PINNs) trained on geometrically nontrivial
PDEs.  Let $\Omega \subset \mathbb{R}^d$ denote the physical domain,
with boundary $\partial\Omega$, and let $u:\Omega\rightarrow\mathbb{R}^m$ be the
solution field.  We consider the PDE problem
of the form
\begin{equation}
    \mathcal{N}[u](x) = f(x),
    \quad x\in\Omega,
    \qquad
    \mathcal{B}[u](x) = g(x),
    \quad x\in\partial\Omega,
    \label{eq:pde-general}
\end{equation}
where $\mathcal{N}$ is the differential PDE operator, $\mathcal{B}$ is the boundary
operator, and $f$ and $g$ are the forcing and boundary terms.
For time-dependent problems, time is treated as an additional input coordinate
and the space-time domain is sampled by the same collocation interface.

\subsection{PINN Objective}
\label{subsec:pinn-objective}

Following the standard PINN formulation \cite{raissi2019physics}, we train a
differentiable parametric model $u_\theta$ with parameters $\theta$ to
approximate the unknown solution field.  The defining feature of a PINN is the
use of automatic differentiation to compute the derivatives required by the
PDE and boundary conditions.  The model is trained by minimizing a composite
loss combining supervised data observations, the interior PDE residual, and
boundary terms:
\begin{equation}
    \mathcal{L}(\theta)
    =
    w_d \mathcal{L}_d(\theta;\mathcal{D})
    + w_r \mathcal{L}_r(\theta;\mathcal{X}_r)
    + w_b \mathcal{L}_b(\theta;\mathcal{X}_b).
    \label{eq:pinn-objective}
\end{equation}
Here $\mathcal{D}$ is a supervised dataset of solution observations, which may
be exact or noisy, $\mathcal{X}_r\subset\Omega$ is the interior collocation set, and
$\mathcal{X}_b\subset\partial\Omega$ is the set used for boundary enforcement.
The interior residual loss evaluates
the PDE constraint violation across the collocation set:
\begin{equation}
    \mathcal{L}_r(\theta;\mathcal{X}_r)
    =
    \frac{1}{|\mathcal{X}_r|}
    \sum_{x\in\mathcal{X}_r}
    \left\|
        \mathcal{R}_\theta(x)
    \right\|_2^2,
    \qquad
    \mathcal{R}_\theta
    =
    \mathcal{N}[u_\theta](x)-f(x).
    \label{eq:residual-loss}
\end{equation}
The central methodological question is how to choose $\mathcal{X}_r$.

\subsection{Mesh-Backed Residual Scoring}
\label{subsec:mesh-backed-adaptive}

Our first method introduces a mesh-backed scaffold for residual-guided
adaptation.  At iteration $k$, the method takes the current PINN $u_{\theta_k}$
and an auxiliary simplicial mesh
$\mathcal{T}^{(k)}=\{E_i^{(k)}\}_{i=1}^{M_k}$, where each element is a triangle
in two spatial dimensions and a tetrahedron in three spatial dimensions, and
converts residual evaluations on that mesh into a spatial importance map.  The
mesh is used only to organize residual information: it provides elementwise
regions, measures, and neighborhoods that let the method reason about coherent
parts of the domain, while leaving the PINN objective in
\cref{eq:pinn-objective} unchanged.

To obtain a stable regional indicator, we evaluate the current residual on a
small set of barycentric quadrature points $Q_i^{(k)}$ inside each element and
define the element score
\begin{equation}
    s_i^{(k)}
    =
    |E_i^{(k)}|^{a}
    \frac{1}{|Q_i^{(k)}|}
    \sum_{q\in Q_i^{(k)}}
    \left\|
        \mathcal{R}_{\theta_k}(q)
    \right\|_2^2,
    \label{eq:element-indicator}
\end{equation}
where $|E_i^{(k)}|$ is area or volume and $a$ controls how strongly the element
measure enters the score.  This aggregation step summarizes local residual
behavior over a region rather than reacting to isolated point evaluations, and
therefore provides a more stable indication of where the current PINN remains
under-resolved.

Because the elementwise scores can still exhibit transient local spikes, we
smooth them over the element adjacency graph:
\begin{equation}
    \tilde{s}_i^{(k)}
    =
    (1-\lambda_s)s_i^{(k)}
    +
    \lambda_s
    \frac{\sum_{j\in\mathcal{N}(i)} |E_j^{(k)}|s_j^{(k)}}
         {\sum_{j\in\mathcal{N}(i)} |E_j^{(k)}|},
    \label{eq:score-smoothing}
\end{equation}
where $\mathcal{N}(i)$ is the set of neighboring elements and
$\lambda_s\in[0,1]$ is the smoothing parameter.  Elements whose smoothed score
exceeds a fixed fraction of the maximum score are marked for refinement.  The
output of this stage is thus a refined, mesh-backed importance map that
identifies coherent high-residual regions and is passed to the subsequent point
allocation stage.

\subsection{Power-Tempered Point Allocation}
\label{subsec:power-tempered}

Our primary adaptive method, denoted \texttt{adaptive\_power\_tempered}, turns
the scaffold scores into a fixed-size collocation set through a rank-persistent
probability distribution.  First, the smoothed element scores are converted into
normalized ranks $z_i^{(k)}\in[0,1]$.  To reduce sensitivity to transient
optimizer noise, ranks are blended with their nearest previous element scores:
\begin{equation}
    \hat{z}_i^{(k)}
    =
    \alpha \hat{z}_{\pi(i)}^{(k-1)}
    +
    (1-\alpha) z_i^{(k)},
    \label{eq:rank-persistent-score}
\end{equation}
where $\alpha\in[0,1]$ is the persistence coefficient and $\pi(i)$ maps a current
element to the nearest element centroid from the previous scaffold.  During the
warmup period, the method uses the current ranks directly.

The persistent ranks define an exponentially tilted element distribution:
\begin{equation}
    p_i^{(k)}
    =
    \frac{
        |E_i^{(k)}|^\gamma
        \exp\!\left(\beta_k \hat{z}_i^{(k)}\right)
    }{
        \sum_{j=1}^{M_k}
        |E_j^{(k)}|^\gamma
        \exp\!\left(\beta_k \hat{z}_j^{(k)}\right)
    },
    \label{eq:apt-distribution}
\end{equation}
where $\gamma$ controls geometric coverage.  The temperature is tied to the
concentration of the raw residual field.  Define
\begin{equation}
    H_k
    =
    -\frac{1}{\log M_k}
    \sum_{i=1}^{M_k}
    \rho_i^{(k)}\log\rho_i^{(k)},
    \qquad
    \rho_i^{(k)}
    =
    \frac{\tilde{s}_i^{(k)}}{\sum_j\tilde{s}_j^{(k)}}.
    \label{eq:residual-entropy}
\end{equation}
We then set
\begin{equation}
    \beta_k
    =
    \beta_{\min}
    +
    \left(\beta_{\max}-\beta_{\min}\right)(1-H_k).
    \label{eq:apt-temperature}
\end{equation}
When residuals are diffuse, the entropy is high and the sampler remains close to
a geometric coverage distribution.  When residuals are concentrated, the
temperature increases and more points are allocated to persistent high-residual
regions.  Some variants mix \cref{eq:apt-distribution} with the true area/volume
distribution as a coverage floor, but the final paper-facing experiments use
the default power-tempered method unless otherwise stated.

Given probabilities $p_i^{(k)}$ and budget $N_r$, integer element quotas are
assigned deterministically by largest remainder:
\begin{equation}
    n_i^{(k)}
    =
    \left\lfloor N_r p_i^{(k)} \right\rfloor
    +
    \delta_i,
    \qquad
    \sum_i n_i^{(k)} = N_r,
    \label{eq:largest-remainder}
\end{equation}
where the $\delta_i$ allocate the remaining points to the elements with the
largest fractional remainders.  Points are then sampled uniformly inside each
selected triangle or tetrahedron using barycentric coordinates.

\subsection{Comparison Samplers}
\label{subsec:comparison-samplers}

We compare the adaptive method against three families of baselines.  Uniform
random sampling draws interior points by rejection sampling from the domain
bounding box.  Halton sampling uses a low-discrepancy sequence with the same
rejection interface and advances to a fresh, non-overlapping sequence window at
each iteration.  RAD follows residual-based adaptive distribution sampling:
\begin{equation}
    \Pr(x_j)
    =
    \frac{
        |\mathcal{R}_{\theta}(x_j)|^\kappa/
        \overline{|\mathcal{R}_{\theta}|^\kappa}
        + c
    }{
        \sum_{\ell}
        \left(
        |\mathcal{R}_{\theta}(x_\ell)|^\kappa/
        \overline{|\mathcal{R}_{\theta}|^\kappa}
        + c
        \right)
    },
    \label{eq:rad-weighting}
\end{equation}
where $\{x_j\}$ is a candidate set, $\kappa$ is the residual exponent,
and $c$ is a coverage regularizer.  In contrast to the mesh-scaffold method,
RAD samples from pointwise residual weights on a candidate cloud rather than
from element-level residual scores.

\subsection{Evaluation Metrics}
\label{subsec:evaluation-metrics}

For each method and seed, the reported checkpoint is selected by the validation
score recorded during training.  We evaluate solution error against the FEM
reference solution on a fixed reference mesh:
\begin{equation}
    e_{\mathrm{rel}}
    =
    \frac{
        \left\|u_\theta-u_{\mathrm{ref}}\right\|_{L^2(\Omega)}
    }{
        \left\|u_{\mathrm{ref}}\right\|_{L^2(\Omega)}
    },
    \qquad
    r_{\mathrm{fixed}}
    =
    \left(
        \frac{1}{|\Omega|}
        \int_{\Omega}
        \left\|\mathcal{R}_{\theta}(x)\right\|_2^2
        \dd x
    \right)^{1/2}.
    \label{eq:evaluation-metrics}
\end{equation}
The fixed residual metric is computed on the same reference discretization for
all methods in a run.  Runtime is recorded per iteration and accumulated through
the selected checkpoint so that sampler overhead is part of the comparison.
